\subsection{Preserve payment history when financing and remedies change}
\label{sec:selected-financing-history}

A company that pays late can end with the same cash, debt and equity as a
company that pays on time. The ending balance sheet does not distinguish the
two histories. A simulator must retain what was due, what remained unpaid
when payment was missed, what was subsequently paid, and which contractual
consequences were actually established. Otherwise a successful funding decision
can silently make a past failure disappear.

Consider the supplied zero-coupon promise used in the financing comparison.
The opening book contains cash of 60, inventory of 100, debt with a carrying
amount of 80 and equity of 80. Contractual face is 100, due on 30 June 2027.
A single explicit measurement interval increases carrying debt and interest
expense by 20 at maturity. Forty new shares are subscribed for 40 in cash.
There are no fees, taxes, penalties or modification effects in this example;
the price and funding date are supplied scenario assumptions. These choices
make payment timing visible without requiring a financing equilibrium.

If the subscription settles on 29 June, the company pays all 100 at maturity.
If it settles on 2 July, the company first pays its available 60 at maturity,
records the unpaid 40 on 1 July, and pays the remaining 40 on 3 July. Both
paths finish with cash and debt equal to zero, inventory of 100 and equity of
100 after the interest expense. Only the second path has a missed-payment
history. The signed journal convention records debt as a credit balance;
the following table displays its positive carrying amount.

\begin{center}
\begin{tabular}{lrrrr}
\toprule
Late-funding path & Cash & Debt & Unpaid claim & Past miss \\
\midrule
Opening & 60 & 80 & 100 & 0 \\
Maturity, before payment & 60 & 100 & 100 & 0 \\
Partial payment, 30 June & 0 & 40 & 40 & 0 \\
Observation, 1 July & 0 & 40 & 40 & 40 \\
Share subscription, 2 July & 40 & 40 & 40 & 40 \\
Final payment, 3 July & 0 & 0 & 0 & 40 \\
\bottomrule
\end{tabular}
\end{center}

The historical amount in the last column is the amount recorded at the
missed-payment observation. It is not the current unpaid balance, and payment
does not reduce it. Neither that amount nor the missed flag is, by itself,
a finding that a particular legal event of default occurred.

\paragraph{Check the state and the change separately.}
The selected financing declaration uses the existing typed requirement
evaluator for 26 snapshot conditions. Face, remaining claim, recorded arrears,
historical missed amount and share counts must be nonnegative. The remaining
claim cannot exceed face. Before maturity the unpaid claim remains face while
the carrying amount may differ; after the selected maturity operation, carrying
debt equals the remaining claim. A recorded miss requires a matured promise,
an observation after its due date and a positive historical missed amount.
The recorded arrears equal the remaining claim in this selected post-observation
state. An unpaid state observed after the due date cannot pass without the
explicit missed-payment observation.

These conditions are useful but insufficient. After full payment, a snapshot
with a historical missed amount of 50 could satisfy the inequalities even
though the recorded amount was 40. Likewise, a nonnegative share count can
remain unchanged when a subscription should have increased it. The transition
declaration therefore binds the old state, typed event inputs and new state.
For a late payment of amount \(p\), let \(U\) be the remaining claim,
\(A\) the recorded arrears and \(H\) the historical missed amount. The
selected change requires
\[
  U_{\mathrm{after}}=U_{\mathrm{before}}-p,\qquad
  A_{\mathrm{after}}=A_{\mathrm{before}}-p,\qquad
  H_{\mathrm{after}}=H_{\mathrm{before}}.
\]
Cash falls by \(p\), the signed debt balance increases by \(p\), and the
missed flag remains true. Original face, due date and old shares are preserved.
A subscription has a different relation: new shares increase by the supplied
share quantity, while cash and capital change by quantity times issue price.
Money and shares have different dimensions, so a money-typed value cannot be
used as a share count merely because both are exact rationals.

The five declared transitions cover maturity, subscription, due-date payment,
missed-payment observation and late payment. They also preserve unrelated
memory, account balances and prior cash-flow records. Deliberately removing
the share-memory or historical-miss write can still yield a balanced proposal
from the generic engine. The selected transition relation rejects both. The
failed proposal remains available for diagnosis, while the accepted current
state and unexecuted command suffix retain the last valid prefix.

\paragraph{An internally consistent history still needs an opening premise.}
Opening assessment binds the selected method registry, schemas, reporting map,
account roles and typed financial state. A separate supplied history record
names the exact opening book and its debt-term fields, with nonblank provenance.
Missing or partial history is unresolved. A paid opening that erases the missed
flag and amount can be internally consistent, but it conflicts with the retained
history record and is refused. A caller who supplies the same false account
in both records can still satisfy this consistency test; source evidence must
establish the historical facts.

A valid intermediate snapshot is also not automatically a fiscal opening.
The existing statement engine requires the opening date to equal the fiscal
period start, an empty current-period cash-flow sequence and zero temporary
account balances. The first verification attempt incorrectly expected an
intermediate funded state to qualify as an opening. Its refusal exposed the
test error. The corrected comparison checks snapshot and supplied-history
consistency separately while preserving the fiscal-opening restrictions.
Closing and carrying the financing history require an explicit further step.

\paragraph{Keep current contractual findings apart from economic history.}
The conditional remedy examples introduce a separate record for an identified
hypothetical series. It contains the selected clause finding, its observation
date, the acceleration finding and the evidence identifier. The applicability
of the selected clause is a supplied premise. A known unpaid principal amount
must agree with the outstanding claim in the economic book; an unknown
applicability finding remains unresolved rather than becoming false.

The retained conditional reference distinguishes seven kinds of question:
clause satisfaction, authorized notice, new acceleration, remaining acceleration
after rescission, protected-change consent, majority waiver and a finding about
another series. Explicit typed predicates reproduce its 58 hypothetical
classifications, including unknown facts and unresolved boundary-day order.
Only three of these routes are composed into state updates here: a clause
observation, a new acceleration assessment and a rescission assessment. The
classification table is not a complete contract execution system.

In the composed example, the funded but still unpaid company first receives
a supplied conditional principal-clause finding. That observation changes no
cash and does not itself accelerate the obligation. A separate, supported
declaration establishes acceleration. The existing typed settlement then pays
the 40. A fresh paid-clause observation becomes unsatisfied, but a lack of a
new acceleration trigger leaves the earlier acceleration finding intact.
The selected rescission assessment changes that state only when all of its
supplied conditions are resolved as required. The historical missed amount
remains 40 throughout. These operations record conditional contractual status;
priority, penalties, accelerated cash amounts and insolvency consequences need
their own economic and accounting methods.

The initial contract-status record and its schema extension are explicit
synthetic premises. They do not establish that the retained issuer description
governs this promise. The source-dependent conditions retain their original
scope and unresolved operative-instrument dependencies. Known automatic and
declaration routes also keep their different information requirements: the
automatic branch does not require a notice finding that it does not use.

\paragraph{A matching date does not establish that a finding is current.}
The first composed candidate required a clause finding from the same date as
the proposed acceleration. Its original trace passed, but a further comparison
found a failure: record the unpaid clause, pay the full claim later in the same
dated sequence, then request acceleration from the old finding. The date still
matches although the economic state has changed. That candidate incorrectly
accepts the stale finding relative to its own current-clause requirement.

The successor records a digest of a fixed economic-state projection with each
resolved clause observation. The projection contains the book's identity and
period, account definitions, balances, cash account, cash-flow history,
non-contract memory and its types. Conditional-status memory and seen-event
bookkeeping are excluded, so recording a remedy does not invalidate itself.
The selected schema and types are checked separately. A new acceleration
requires both the date and economic-state digest to match the stored observation.

The same-day payment now makes the stale request unresolved before any accepted
write. A new observation of the paid claim supports the consistent assessment
that there is no new trigger. Balanced account changes, changes to unrelated
economic memory and revisions to cash-flow evidence also invalidate the old
observation. This conservative rule avoids guessing which economic changes
matter. Matching digests establish equality of the supplied record, not truth
of the evidence or persistence of external legal facts. The relevant external
findings still have to be supplied for the event.

The selected financing contract reproduces four complete financing paths on
two account charts, including their statements. The repaired remedy composition
reproduces eight complete transitions, with separate same-day failure and repair
evidence. These checks support the stated finite declarations and preserve the
rejected predecessor. Wider financing decisions, carry, changed rights, actual
series admission and independent review remain necessary for the full design.

\begingroup\raggedright
The independent settlement and conditional-source records are retained under
\path{data/accounting-master/gap-program-20261026/}.
The selected financial conditions and opening checks are under
\path{data/accounting-master/financing-opening-invariants-20261026/}.
The date-only failure and its preserved comparator are under
\path{data/accounting-master/financing-remedy-binding-20261026/};
the checked successor is under
\path{data/accounting-master/financing-remedy-state-binding-20261026/}.
\par\endgroup
